\documentclass[10pt,a4paper]{amsart}
\usepackage[margin=19mm]{geometry}
\usepackage{amsmath,amssymb,mathtools}
\usepackage[scr=rsfs,cal=euler]{mathalfa}
\usepackage{xfrac}
\usepackage[unicode,hidelinks,bookmarks=false]{hyperref}
\newcommand{\ie}{i.e.\ }
\newcommand{\cf}{cf.\ }
\newcommand{\resp}{respectively\ }
\newcommand{\Subsection}{\subsection}
\providecommand{\nobreakdash}{\nobreak\hbox{-}}
\setlength{\emergencystretch}{2em}
\multlinegap0pt
\usepackage{bm}
\usepackage{bbm}
\usepackage{dsfont}
\usepackage{xspace}
\usepackage{cancel}
\usepackage{mathtools}
\usepackage{amsthm}
\makeatletter
\@ifundefined{lemma}{\newtheorem{lemma}{Lemma}}{}
\@ifundefined{theorem}{\newtheorem{theorem}{Theorem}}{}
\makeatother
\newcommand{\R}{\mathbb{R}}
\newcommand{\be}{\mathbf e}
\newcommand{\bu}{\mathbf u}
\newcommand{\bv}{\mathbf v}
\title[Discretization error analysis for a radially symmetric harmo]{Discretization error analysis for a radially symmetric harmonic map heat flow problem}
\author{Nam Anh Nguyen, Arnold Reusken}
\date{Source: arXiv:2503.03525v2 (CC BY 4.0)}
\begin{document}
\maketitle

\noindent\emph{Source and licence.} Discretization error analysis for a radially symmetric harmonic map heat flow problem. Version arXiv:2503.03525v2 (CC BY 4.0), distributed under the Creative Commons Attribution 4.0 International licence, as recorded in the arXiv licence metadata for this exact version and shown on the abstract page https://arxiv.org/abs/2503.03525v2. Full rights notice: licenses/arxiv-2503.03525-CC-BY-4.0.txt.\par\smallskip

\noindent\textit{Editorial scope.} Counted unit: the Theorem whose source label is \texttt{mainthm} in the article (source0.tex lines 395--402, source label \texttt{mainthm}), delivered together with the article's own complete original proof of it (source0.tex lines 403--488, one \texttt{proof} environment of 86 lines). No mathematics is written, repaired or re-derived: statement and proof are the article's own text line for line. Required context retained uncounted: lemma4 (lines 322--327); lemma5 (lines 341--346); lemma8 (lines 358--363); erreq2 (lines 388--391). Every definition, display and label of those lines is retained unchanged. Omitted from this package and not redistributed: 1--321, 328--340, 347--357, 364--387, 392--394, 489--702. Nothing inside a retained excerpt is omitted or altered: every retained body is delivered exactly as the article's lines are written, so no omission receipt of any kind accompanies this package. Citations: the retained excerpts carry no citation command at all, so no bibliography entry of the article is transcribed and no reference is redistributed. Reference adaptation: none was needed and none was made; every reference inside the delivered unit resolves inside it, so no unresolved reference or citation is produced. Printed slips kept literal: no printed word, symbol, hypothesis or number of the counted statement or of its proof was altered. Layout and class: the article's own class is \texttt{siamltex} with packages including amsmath, amsfonts, amssymb, latexsym, graphicx, overpic, subfigure, caption, pgfplots, color, booktabs, tikz; the wrapper is the lane's pinned \texttt{amsart} editorial wrapper at 10pt on A4, which restates the theorem-like environments the retained text needs and adds the article's own macro definitions that the retained text needs (\texttt{\textbackslash R}, \texttt{\textbackslash be}, \texttt{\textbackslash bu}, \texttt{\textbackslash bv}), with extra wrapper packages (bm, bbm, dsfont, xspace, cancel, mathtools). The delivered numbering is the unit's own. Assets: no retained span contains a figure environment, a graphics-inclusion command, a PDF-page inclusion, a table or a TikZ picture; no image file is copied into this package, the package declares no asset dependency and no image file is redistributed with it. Licence: version arXiv:2503.03525v2 of this article is distributed under the Creative Commons Attribution 4.0 International licence, as recorded in the arXiv licence metadata for this exact version and shown on the abstract page https://arxiv.org/abs/2503.03525v2. A full rights notice is delivered at licenses/arxiv-2503.03525-CC-BY-4.0.txt. Characters: every retained excerpt is pure ASCII, so the delivered engine, XeLaTeX with its default Latin Modern text font, reports no missing character.
\begin{lemma} \label{lemma4}
 Take $\alpha \in [0,1)$. Let $c_\alpha \in (0,\tfrac{\pi}{2}]$ be such that $\frac{\sin(2c_\alpha)}{2c_\alpha}=\alpha^2$.  Assume $\bu^0$ satisfies $\|\bu^0\|_\infty \leq c_\alpha$. Then the following holds:
 \[
   \|\bu_\alpha^n\|_\infty \leq \|\bu_\alpha^0\|_\infty \quad \text{for all}~n \geq 0.
 \]
\end{lemma}

 \begin{lemma} \label{lemma5}
  The following holds
  \begin{equation} \label{bound1}
    \|\big(I+ \Delta t(D^\frac12 C D^{- \frac12} +G(\bu^n)D^{-2})\big)^{-1}\|_2 \leq 1.
  \end{equation}
 \end{lemma}

\begin{lemma} \label{lemma8}
 Define $g(y):=\frac{\sin(2y)}{2y}$, $y \neq 0$, $g(0)=1$. The following holds:
 \[ 
  |g(y) - g(z)| \leq \tfrac43 \max \{|y|,|z|\} |y-z| \quad \text{for all}~ y,z \in \R.
 \]
\end{lemma}

\begin{equation} \label{erreq2} \begin{split}
 \frac{\be^{n+1}- \be^n}{\Delta t}  & = - C\be^{n+1} - \left(G(\bu(t_{n+1}))D^{-2}\bu(t_{n+1})- G(\bu^n)D^{-2}\bu^{n+1}\right)  \\  & \quad  -\Delta t \,\be_{\partial t}(t_{n+1}) +   h^2 \be_{\partial^2 x}(t_{n+1})+h^2 D^{-1}\be_{\partial x}(t_{n+1}), \quad 0 \leq n \leq M-1.
 \end{split}
\end{equation}

\begin{theorem} \label{mainthm}    Take $\alpha  \in (\tfrac12, 1)$ and assume $\|\bu^0\|_\infty \leq c_\alpha$, with $c_\alpha$ as in Lemma~\ref{lemma4}.
Define $d_\alpha:=\|D^{-\alpha}\bu^0\|_\infty$.  
For the error $\be^n$ the following holds:
\begin{equation} \label{optboound}
  \|\be^n\|_{D,h} \leq   c \big( \Delta t + h^2 |\ln h|^\frac12), \quad 1 \leq n \leq M,
\end{equation}
with a constant $c$ that depends only on $\|\frac{\partial^j u}{\partial  x^j}\|_{L^\infty(I\times[0,T])}$,  $0 \leq j \leq 4$, $\| \frac{\partial^j u}{\partial t^j}\|_{L^\infty(I\times[0,T])}$,  $0 \leq j \leq 2$,  $T$, $\alpha$, $c_\alpha$ and $d_\alpha$. The dependence of $c$ on these quantities can be deduced from the proof.
\end{theorem}

\begin{proof} We use $c$ for a (varying) constant that depends only on the quantities mentioned above. 
From \eqref{erreq2} we obtain, after multiplication with $D^\frac12$
\[ \begin{split}
  & (I+\Delta t D^\frac12C D^{-\frac12})D^\frac12\be^{n+1} \\ & =D^{\frac12}\be^n - \Delta t\, E    -(\Delta t)^2 \,D^\frac12 \be_{\partial t}(t_{n+1}) +  \Delta t\,  h^2 D^\frac12\be_{\partial^2 x}(t_{n+1})+ \Delta t\, h^2 D^{-\frac12}\be_{\partial x}(t_{n+1}), \\
  & E:= \left(G(\bu(t_{n+1}))D^{-\frac32}\bu(t_{n+1})- G(\bu^n)D^{-\frac32}\bu^{n+1}\right).
\end{split} \]
For linearization of the term $E$ we use the splitting
\begin{equation} \label{splitE}
 \begin{split}
  E & = G(\bu(t_{n+1}))D^{-\frac32}\bu(t_{n+1})- G(\bu(t_n))D^{-\frac32}\bu(t_{n+1}) \\
     & + G(\bu(t_n))D^{-\frac32}\bu(t_{n+1})- G(\bu^n)D^{-\frac32}\bu(t_{n+1}) \\
     & + G(\bu^n)D^{-\frac32}\bu(t_{n+1})- G(\bu^n)D^{-\frac32}\bu^{n+1} =:E_1+E_2+E_3.
 \end{split}
\end{equation}
The term $E_3$ is shifted to the left hand side. We thus get  
\begin{equation} \label{erreq8} \begin{split}
  & \big(I+\Delta t (D^\frac12C D^{-\frac12} +G(\bu^n)D^{-2})\big)D^{\frac12}\be^{n+1} \\ & =D^{\frac12}\be^n - \Delta t\, (E_1+E_2)  \\ &   -(\Delta t)^2 \,D^\frac12 \be_{\partial t}(t_{n+1}) +  \Delta t\,  h^2 D^\frac12\be_{\partial^2 x}(t_{n+1})+ \Delta t\, h^2 D^{-\frac12}\be_{\partial x}(t_{n+1}).
\end{split} \end{equation}
We  consider the term $E_1$. Note that the $\|\cdot\|_{2,h}$ matrix norm is scaling invariant. Replacing it by $\|\cdot\|_2$ we obtain
\[
 \|E_1\|_{2,h} \leq \|\big(G(\bu(t_{n+1}))- G(\bu(t_n))\big)D^{-\frac12}\|_2 \|D^{-1}\bu(t_{n+1})\|_{2,h}.
\]
Note that $\|\bv\|_{D,h}\leq \|\bv\|_{2,h} \leq \|\bv\|_\infty$ holds for all $\bv \in \R^N$. 
For the continuous solution $\bu(t_n)$ we have
\begin{equation} \label{erreq4} \begin{split}
  \|D^{-1}\bu(t_n)\|_{\infty} & = \max_{1 \leq i \leq N} \left| \frac{u(x_i,t_n)}{x_i}\right| \\ & =\max_{1 \leq i \leq N} \frac{1}{x_i} \left|
  \int_0^{x_i} u_x(s,t_n)\, ds\right| \leq \|\frac{\partial u}{\partial  x}\|_{L^\infty(I\times[0,T])}=c_1.
\end{split} \end{equation}
Using $x_i^{-\frac12} \leq x_i^{-1}$,  Lemma~\ref{lemma8} and  \eqref{erreq4} we obtain
\[ \begin{split}
 & \|\big(G(\bu(t_{n+1})) - G(\bu(t_{n}))\big)D^{-\frac12}\|_2  \leq  \max_{1 \leq i \leq N} \frac{1}{x_i} |g(u(x_i,t_{n+1}))-g(u(x_i,t_{n}))| \\ &  \leq \tfrac43 c_1 \max_{1 \leq i \leq N} |u(x_i,t_{n+1})-u(x_i,t_n)| \\ & \leq \tfrac43 c_1 \Delta t \|u_t\|_{L^\infty(I \times [0,T])} = c \Delta t.
\end{split}\]
Using this and \eqref{erreq4} yields
\begin{equation} \label{erreq5}
 \|E_1\|_{2,h} \leq c\, \Delta t.
\end{equation}
For the term $E_2$ we note, for $1 \leq i \leq N$,
\[ \begin{split}
 |(E_2)_i| & =\left| \big(g(u(x_i,t_{n})) - g(u_i^n)\big)\big)x_i^{-\frac32} u(x_i,t_{n+1})\right| \\
 & \leq \tfrac43 \max \{|u(x_i,t_{n})|,|u_i^n|\} |e_i^n| x_i^{-\frac32} |u(x_i,t_{n+1})| 
  \\ &  =\tfrac43 \max \{|u(x_i,t_{n})|,|u_i^n|\} x_i^{-\alpha} |e_i^n| x_i^{\alpha-\frac32} |u(x_i,t_{n+1})| 
\end{split}\]
Using $x_i^{-\alpha}|u(x_i,t_{n})| \leq x_i^{-1}|u(x_i,t_{n})| \leq  c_1$ and 
\begin{equation} \label{newlabel} x_i^{-\alpha} |u_i^n| \leq \|D^{-\alpha} \bu^n\|_\infty \leq \|D^{-\alpha} \bu^0\|_\infty= d_\alpha,
\end{equation}
cf. Lemma~\ref{lemma4},  this yields
\[
 |(E_2)_i| \leq \tfrac43 \max\{c_1, d_\alpha\} |x_i^\frac12 e_i^n| |x_i^{\alpha-2}u(x_i,t_{n+1})|.
\]
With Cauchy-Schwarz we thus get
\begin{equation} \label{estE2}
 \|E_2\|_{2,h} \leq  \tfrac43 \max\{c_1, d_\alpha\} \|\be^n\|_{D,h} \left(h \sum_{i=1}^N x_i^{2\alpha-4} u(x_i,t_{n+1})^2 \right)^\frac12.
\end{equation}
Using $ \max_{i}|\frac{u(x_i,t_{n+1})}{x_i}| \leq c_1$ and $x_i=ih$  we obtain for the expression above
\begin{align*} 
  & \left(h \sum_{i=1}^N x_i^{2\alpha-4} u(x_i,t_{n+1})^2 \right)^\frac12 \leq c_1 \left(h \sum_{i=1}^N x_i^{2(\alpha-1)}  \right)^\frac12=c_1 \left(h^{2 \alpha -1} \sum_{i=1}^N i^{2(\alpha-1)}  \right)^\frac12 \nonumber \\
 & \leq c_1 \left(h^{2 \alpha -1}\big( 1+ \int_1^{N} z^{2 (\alpha-1)}\, dz \big) \right)^\frac12 \leq \frac{c_1}{\sqrt{2 \alpha -1}}.
\end{align*}
Summarizin*g we obtain
\begin{equation} \label{R5}
   \|E_2\|_{2,h} \leq c \|\be^n\|_{D,h}, \quad \text{with}~ c= \tfrac43 \max\{c_1, d_\alpha\}\frac{c_1}{\sqrt{2 \alpha -1}}.
\end{equation}
Using  the results \eqref{erreq5} and \eqref{R5} in \eqref{erreq8} and with the stability estimate of Lemma~\ref{lemma5} we thus get  
\begin{equation}\label{eq7} \begin{split}
 & \|\be^{n+1}\|_{D,h} \leq (1 + c \Delta t)\|\be^n\|_{D,h} + c\, (\Delta t)^2 \\
  & + (\Delta t)^2 \|D^\frac12 \be_{\partial t}(t_{n+1})\|_{2,h} +  \Delta t\, h^2\| D^\frac12\be_{\partial^2 x}(t_{n+1})\|_{2,h}+ \Delta t\, h^2 \|D^{-\frac12}\be_{\partial x}(t_{n+1})\|_{2,h}.
\end{split} \end{equation}
We finally estimate the three terms in the second line of \eqref{eq7}. We have
\[ \|D^\frac12 \be_{\partial t}(t_{n+1})\|_{2,h} \leq \|\be_{\partial t}(t_{n+1})\|_\infty \leq \frac12 \| \frac{\partial^2 u}{\partial t^2}\|_{L^\infty(I\times[0,T])},
\]
and  
\[
 \| D^\frac12\be_{\partial^2 x}(t_{n+1})\|_{2,h} \leq \| \be_{\partial^2 x}(t_{n+1})\|_\infty \leq \frac{1}{12}\|\frac{\partial^4 u}{\partial  x^4}\|_{L^\infty(I\times[0,T])}.
\]
In  the third term there is a scaling with $D^{-\frac12}$. For this term we obtain, using $\sum_{i=1}^N \frac{1}{i} \leq 1 +\int_1^N \frac{1}{x} \, dx = 1 + \ln N \leq 1 + |\ln h| \leq c |\ln h|$:
\[ \begin{split} 
  \|D^{-\frac12}\be_{\partial x}(t_{n+1})\|_{2,h} & = \frac16 \left( h \sum_{i=1}^N \frac{1}{x_i} \frac{\partial^3 u}{\partial x^3}(\xi_i,t_{n+1})^2 \right)^\frac12 
   \leq \frac{1}{6}\|\frac{\partial^3 u}{\partial  x^3}\|_{L^\infty(I\times[0,T])} \left( \sum_{i=1}^N \frac{1}{i}\right)^\frac12 \\
  & \leq c \|\frac{\partial^3 u}{\partial  x^3}\|_{L^\infty(I\times[0,T])} |\ln h|^\frac12.
\end{split} \]
Using these estimates in \eqref{eq7} we obtain
\[
  \|\be^{n+1}\|_{D,h} \leq (1 + c \Delta t)\|\be^n\|_{D,h} + c\, \Delta t \big( \Delta t + h^2 |\ln h|^\frac12\big).
\]
A standard recursive argument, using $(1+ c \Delta t)^M \leq e^{c T}$ and $\be^0=0$ completes the proof.
\end{proof}

\end{document}
